\FloatBarrier
\section{تفسیرپذیری در \lr{v2}}
\label{sec:interp}

دقت به ما می‌گوید مدل چقدر خوب پیش‌بینی می‌کند. تفسیرپذیری می‌پرسد مدل برای این پیش‌بینی به کدام ورودی‌ها تکیه کرده است: کدام لحظه‌ها، کدام الکترودها، کدام باندها و کدام عضله‌ها. این بخش روش‌ها را توضیح می‌دهد، نشان می‌دهد چطور خودِ روش‌ها سنجیده شدند و نتایج را برای همه‌ی شرکت‌کننده‌ها و با جزئیات برای \lr{P1} و \lr{P2} می‌آورد.

\subsection{چه چیزی توضیح داده می‌شود}

هر توضیح درباره‌ی یک عدد مشخص است: خروجی $o$ (یکی از $K$، $E$ یا $F_{eeg}$)، برای عضله‌ی $c$، در \lr{token} $t$ از یک پنجره‌ی ۸ ثانیه‌ای:
\begin{equation}
y=\mathrm{out}_o[t,c].
\end{equation}
پنجره‌ها دور رخدادهای مهم آزمایش ساخته می‌شوند: «اولین تماس» (انگشت به جسم می‌رسد) و «جدا شدن» (\lr{lift-off}، جسم از سطح بلند می‌شود). $t$ وسط پنجره است، یعنی خودِ لحظه‌ی رخداد. پس پرسش این است: «مدل برای پیش‌بینی فعالیت عضله در لحظه‌ی رخداد، به کجای ورودی نگاه کرده است؟»

مسیری که هر خروجی از ورودی طی می‌کند:
\begin{itemize}
\item $K$: \lr{stem}، سه لایه‌ی زمانی سینماتیک، دو لایه‌ی \lr{muscle graph}؛
\item $E$: دو لایه‌ی \lr{electrode graph}، \lr{attention pooling}، دو لایه‌ی زمانی \lr{EEG}، دو لایه‌ی \lr{muscle graph}؛
\item $F_{eeg}$: \lr{electrode graph}، \lr{attention pooling}، لایه‌های زمانی \lr{EEG} و \lr{cross-attention}.
\end{itemize}
\lr{recorder} در یک بار اجرای مدل، ماتریس توجه همه‌ی این لایه‌ها را برمی‌دارد: $A^{h}_{l}\in\mathbb{R}^{n\times n}$ برای لایه‌ی $l$ و سر $h$. همه‌ی روش‌های زیر از همین ماتریس‌ها یا از گرادیان $y$ نسبت به آن‌ها ساخته می‌شوند. شکل~\ref{fig:interp-pipeline} کل فرایند را نشان می‌دهد.

\begin{figure}[htbp]
\centering
% Interpretability pipeline of v2.
\begin{latin}
\begin{tikzpicture}[node distance=5mm and 3mm, font=\footnotesize]
\node[io, text width=6cm] (ev) {test series and their lift events};
\node[box, below=of ev, text width=6cm] (win) {one 8 s window centred on each event};
\node[box, below=of win, text width=6cm] (rec) {one forward pass with the recorder:\\ $A^{h}_{l}$, $\partial y/\partial A^{h}_{l}$ and $\partial y/\partial x$};
\node[box, below=10mm of rec, text width=2.4cm] (v3) {attention flow};
\node[box, left=of v3, text width=2.4cm] (v2) {gradient-weighted rollout};
\node[box, left=of v2, text width=2.4cm] (v1) {rollout};
\node[box, right=of v3, text width=2.4cm] (v4) {gradient $\times$ input};
\node[box, right=of v4, text width=2.4cm] (v5) {per-layer view};
\node[box, below=10mm of v3, text width=6cm] (avg) {mean and SEM over events};
\node[box, below=of avg, text width=6cm] (chk) {deletion check and weight-shuffle check};
\node[box, below=of chk, text width=6cm] (con) {event against rest contrast (E only)};
\node[result, below=of con, text width=6cm] (out) {tables and figures per participant};
\draw[arr] (ev)--(win); \draw[arr] (win)--(rec);
\foreach \v in {v1,v2,v3,v4,v5}{\draw[arr] (rec)--(\v); \draw[arr] (\v)--(avg);}
\draw[arr] (avg)--(chk); \draw[arr] (chk)--(con); \draw[arr] (con)--(out);
\end{tikzpicture}
\end{latin}

\caption{فرایند تفسیر در \lr{v2}: هر روش از یک اجرای ثبت‌شده ساخته می‌شود، روی رخدادها میانگین گرفته می‌شود و پیش از هر ادعا، خود روش سنجیده می‌شود.}
\label{fig:interp-pipeline}
\end{figure}

\subsection{\lr{attention rollout} و \lr{attention flow}}

این دو روش توجه همه‌ی لایه‌ها را پشت سر هم ترکیب می‌کنند تا از خروجی به ورودی برسند. چون پایه‌ی بیشتر نتایج این گزارش‌اند، بخش جداگانه‌ی~\ref{sec:rollout-flow} آن‌ها را با جزئیات توضیح می‌دهد. برای این بخش همین کافی است: \lr{rollout} در هر لایه ماتریس $\hat A_l=\tfrac12\overline{A}_l+\tfrac12 I$ را می‌سازد (میانگین سرها، به‌علاوه‌ی راه مستقیم \lr{residual connection}) و این ماتریس‌ها را در هم ضرب می‌کند: $R=\hat A_L\cdots\hat A_1$. نتیجه برای هر \lr{token} ورودی یک سهم بین ۰ و ۱ است و جمع سهم‌ها ۱ می‌شود. این سهم برای هر پنج عضله‌ی یک جریان یکسان است، چون توجه زمانی بین عضله‌ها مشترک است.

\subsection{\lr{gradient-weighted rollout}}

\textbf{ایده.} \lr{plain rollout} همه‌ی توجه‌ها را به یک اندازه مهم می‌داند. اما ممکن است یک سر توجه، اطلاعاتی را جابه‌جا کند که برای عضله‌ی موردنظر اصلاً اهمیتی ندارد. \lr{gradient-weighted rollout} \cite{chefer2021} این را اصلاح می‌کند: هر وزن توجه در گرادیان خروجی نسبت به همان وزن ضرب می‌شود. گرادیان می‌گوید اگر این وزن کمی بیشتر شود، خروجی چقدر و در کدام جهت تغییر می‌کند. فقط بخش مثبت نگه داشته می‌شود، یعنی توجه‌هایی که خروجی را بالا می‌برند. چون گرادیان برای هر عضله فرق می‌کند، نقشه هم برای هر عضله جداست.

\textbf{ریاضی.}
\begin{equation}
G_l=\frac1H\sum_h\Big(A^{h}_l\odot\frac{\partial y}{\partial A^{h}_l}\Big)^{+},\qquad
\tilde A_l=\frac{G_l}{G_l\mathbf 1},\qquad
R^{\nabla}=\prod_{l=L}^{1}\Big(\tfrac12\tilde A_l+\tfrac12 I\Big).
\end{equation}
$(\cdot)^{+}$ مقادیر منفی را صفر می‌کند. تقسیم بر $G_l\mathbf 1$ هر سطر را دوباره به جمع ۱ می‌رساند؛ سطری که همه‌اش صفر شده همان $I$ می‌ماند. بقیه‌ی کار دقیقاً مثل \lr{rollout} است. گرادیان‌ها با \lr{recorder} در حالت \lr{grad=True} گرفته می‌شوند.

\begin{figure}[htbp]
\centering
% Gradient-weighted rollout.
\begin{latin}
\begin{tikzpicture}[node distance=4mm and 4mm, font=\footnotesize, box/.append style={text width=2.45cm}]
\node[box] (a) {$A^{h}_l$ and $\partial y/\partial A^{h}_l$};
\node[box, right=of a] (b) {element-wise product, positive part};
\node[box, right=of b] (c) {mean over heads, rows rescaled to sum 1};
\node[box, right=of c] (d) {$\tfrac12\tilde A_l+\tfrac12I$, product over layers};
\node[result, right=of d] (e) {map for one muscle};
\draw[arr] (a)--(b); \draw[arr] (b)--(c); \draw[arr] (c)--(d); \draw[arr] (d)--(e);
\end{tikzpicture}
\end{latin}

\caption{\lr{gradient-weighted rollout}.}
\end{figure}

\subsection{نمای هر لایه}

\lr{rollout} لایه‌ها را با هم ترکیب می‌کند و نمی‌گوید کدام لایه چه کرده است. نمای هر لایه همین را نشان می‌دهد: توجه هر لایه و هر سر در لحظه‌ی رخداد، بدون ترکیب. برای لایه‌های زمانی سطر $t$ از $A^{h}_l$، و برای \lr{electrode graph} و \lr{muscle graph} کل ماتریس رسم می‌شود.

\subsection{\lr{gradient × input}}

\textbf{ایده.} \lr{rollout} فقط لایه‌های توجه را می‌بیند. لایه‌هایی مثل \lr{embedding}، \lr{feed-forward} و \lr{decoder} در آن غایب‌اند، و \lr{rollout} نمی‌تواند بگوید کدام ویژگی ورودی (مثلاً کدام باند \lr{EEG}) مهم بوده است. \lr{gradient × input} \cite{shrikumar2016} از کل مدل استفاده می‌کند: مقدار هر ورودی را در گرادیان خروجی نسبت به آن ضرب می‌کند. نتیجه تقریباً می‌گوید اگر این ورودی صفر شود، خروجی چقدر تغییر می‌کند.

\textbf{ریاضی.} برای ورودی $x$ (سینماتیک، تاریخچه یا \lr{band power} \lr{EEG}):
\begin{equation}
s_i=\Big|x_i\,\frac{\partial y}{\partial x_i}\Big|.
\end{equation}
برای هر ویژگی سینماتیک، $s$ روی زمان جمع زده می‌شود. برای هر باند \lr{EEG}، روی زمان و همه‌ی الکترودها. سپس عددها به سهم (با جمع ۱) تبدیل می‌شوند.

\begin{figure}[htbp]
\centering
% Gradient x input.
\begin{latin}
\begin{tikzpicture}[node distance=4mm and 4mm, font=\footnotesize, box/.append style={text width=2.45cm}]
\node[box] (a) {input $x$ (kinematics, history, band power)};
\node[box, right=of a] (b) {back-propagate from $y$: $\partial y/\partial x$};
\node[box, right=of b] (c) {$|x\odot\partial y/\partial x|$};
\node[box, right=of c] (d) {sum over time (and electrodes)};
\node[result, right=of d] (e) {share of each feature or band};
\draw[arr] (a)--(b); \draw[arr] (b)--(c); \draw[arr] (c)--(d); \draw[arr] (d)--(e);
\end{tikzpicture}
\end{latin}

\caption{\lr{gradient × input}.}
\end{figure}

\subsection{میانگین روی رخدادها و مقایسه با استراحت}

\textbf{میانگین روی رخدادها.} نقشه‌ی یک پنجره‌ی تنها پرنویز است، پس هیچ نقشه‌ای از یک پنجره گزارش نمی‌شود. برای هر رخداد یک پنجره ساخته می‌شود، نقشه‌اش به سهم (جمع ۱) تبدیل می‌شود و سپس میانگین و خطای استاندارد ($\mathrm{SEM}=\sigma/\sqrt n$) روی همه‌ی رخدادها گرفته می‌شود. رخدادی که پنجره‌ی ۸ ثانیه‌ای‌اش در سری جا نشود کنار گذاشته می‌شود.

\textbf{مقایسه با استراحت (فقط برای \lr{E}).} حتی اگر یک الکترود سهم بالایی داشته باشد، شاید همیشه سهم بالایی داشته باشد و ربطی به حرکت نداشته باشد. برای جدا کردن این دو، هر بلندکردن با یک پنجره‌ی استراحت بعد از آن (وسط فاصله‌ی دو بلندکردن) جفت می‌شود. برای هر الکترود $e$ تفاوت سهم در رخداد و استراحت ($d_{ie}=a_{ie}-b_{ie}$) و آماره‌ی $t_e=\bar d_e/(s_e/\sqrt n)$ حساب می‌شود. $t$ بزرگ یعنی این الکترود در رخداد به‌طور پایدار سهم بیشتری دارد تا در استراحت.

برای $p$، علامت تفاوت هر جفت ۱۰۰۰ بار به‌طور تصادفی برگردانده می‌شود. هر بار بزرگ‌ترین $|t|$ روی همه‌ی ۳۲ الکترود نگه داشته می‌شود و $p$ هر الکترود از مقایسه با این بیشینه‌ها به دست می‌آید (\lr{max-T}، \cite{nichols2002}). این کار دو مزیت دارد. اول، خطای ناشی از آزمودن ۳۲ الکترود با هم را کنترل می‌کند. دوم، چون رخداد و استراحت از یک سری و یک مدل‌اند، هر اثری که در هر دو مشترک است (مثلاً ساختار خود پنجره) در تفاوت حذف می‌شود.

\begin{figure}[htbp]
\centering
% Event against rest contrast.
\begin{latin}
\begin{tikzpicture}[node distance=4mm and 4mm, font=\footnotesize, box/.append style={text width=2.45cm}]
\node[box] (a) {electrode shares at the event and at the rest after it};
\node[box, right=of a] (b) {paired difference $d_{ie}$};
\node[box, right=of b] (c) {$t_e$ per electrode};
\node[box, right=of c] (d) {1000 random sign flips, keep $\max_e |t|$};
\node[result, right=of d] (e) {corrected $p$ per electrode};
\draw[arr] (a)--(b); \draw[arr] (b)--(c); \draw[arr] (c)--(d); \draw[arr] (d)--(e);
\end{tikzpicture}
\end{latin}

\caption{مقایسه‌ی رخداد با استراحت.}
\end{figure}

\subsection{سنجش خودِ روش‌ها}

هر روش تفسیری یک نقشه می‌دهد، ولی این به معنای درست بودن نقشه نیست. دو آزمون می‌سنجند که نقشه‌ها واقعاً چیزی درباره‌ی مدل آموزش‌دیده می‌گویند.

\textbf{\lr{deletion check}} \cite{samek2017}. اگر یک روش درست بگوید کدام ورودی‌ها مهم‌اند، حذف همان ورودی‌ها باید خروجی را بیشتر از حذف ورودی‌های تصادفی تغییر دهد. «حذف» یعنی صفر کردن، که چون داده‌ها استانداردند برابر میانگین آموزش است. برای \lr{EEG} هر خانه یک الکترود در یک \lr{token} است (هر چهار باند با هم)، و برای سینماتیک همه‌ی ویژگی‌های یک \lr{token}. ۵، ۱۰، ۲۰ و ۴۰ درصد از مهم‌ترین خانه‌ها حذف و با حذف همان تعداد خانه‌ی تصادفی (میانگین ۱۰ بار) مقایسه می‌شوند:
\begin{equation}
\Delta_{top}(q)=\big|\,y(x)-y(x_{\setminus \mathrm{top}(q)})\,\big|,\qquad \Delta_{rand}(q)=\operatorname{mean}_{10}\big|\,y(x)-y(x_{\setminus \mathrm{rand}(q)})\,\big|.
\end{equation}
$x_{\setminus \mathrm{top}(q)}$ ورودی بدون $q$ درصد خانه‌های برتر است و $x_{\setminus \mathrm{rand}(q)}$ ورودی بدون همان تعداد خانه‌ی تصادفی. هر چه $\Delta_{top}$ یک روش بزرگ‌تر باشد، آن روش ورودی‌های واقعاً اثرگذار را بهتر پیدا کرده است.

\textbf{\lr{weight-shuffle check}} \cite{adebayo2018}. یک نسخه از مدل ساخته می‌شود که درایه‌های هر ماتریس وزن آن به‌هم ریخته شده است؛ یعنی مدلی با همان معماری که هیچ چیز یاد نگرفته. نقشه‌ها برای این مدل هم ساخته و با نقشه‌های مدل واقعی مقایسه می‌شوند (\lr{Spearman correlation}). اگر نقشه‌ی مدل به‌هم‌ریخته هم همان شکل را داشته باشد، آن شکل از معماری آمده و نه از آنچه مدل یاد گرفته؛ پس درباره‌ی داده چیزی نمی‌گوید. همبستگی خوب نزدیک صفر است.

\begin{figure}[htbp]
\centering
% Deletion check (top) and weight-shuffle check (bottom).
\begin{latin}
\begin{tikzpicture}[node distance=4mm and 4mm, font=\footnotesize, box/.append style={text width=2.45cm}]
\node[box] (a) {map of one method};
\node[box, right=of a] (b) {rank the input cells};
\node[box, right=of b] (c) {set the top $q$\,\% and a random $q$\,\% to 0};
\node[box, right=of c] (d) {$\Delta_{top}$ and $\Delta_{rand}$};
\node[result, right=of d] (e) {deletion curve over $q$};
\draw[arr] (a)--(b); \draw[arr] (b)--(c); \draw[arr] (c)--(d); \draw[arr] (d)--(e);
\node[box, below=9mm of b] (f) {shuffle every weight matrix};
\node[box, right=of f] (g) {build the map again};
\node[result, right=of g] (h) {Spearman $\rho$ with the original map};
\draw[arr] (a.south)|-(f.west); \draw[arr] (f)--(g); \draw[arr] (g)--(h);
\end{tikzpicture}
\end{latin}

\caption{\lr{deletion check} (بالا) و \lr{weight-shuffle check} (پایین).}
\end{figure}

\subsection{نتایج برای همه‌ی شرکت‌کننده‌ها}

\paragraph{کدام روش ورودی‌های مهم را بهتر پیدا می‌کند؟} جدول~\ref{tab:interp-del} پاسخ \lr{deletion check} است. برای \lr{E}، \lr{gradient-weighted rollout} با فاصله‌ی زیاد بهترین روش است: حذف خانه‌های برتر آن خروجی را \lr{۰٫۲۰۵} تغییر می‌دهد، در حالی که این عدد برای روش‌های دیگر حدود \lr{۰٫۰۹} و برای حذف تصادفی \lr{۰٫۰۲۱} است. این روش در هر ۱۲ شرکت‌کننده اول است (جدول~\ref{tab:interp-parts}). برای \lr{K}، سه روش \lr{rollout}، \lr{gradient-weighted rollout} و \lr{gradient × input} تقریباً برابرند و \lr{attention flow} ضعیف‌تر است. همه‌ی روش‌ها از حذف تصادفی بهترند.

\begin{table}[htbp]
\centering
\small
\caption{\lr{deletion check}، میانگین روی ۱۲ شرکت‌کننده، پنج عضله و چهار کسر حذف (۵، ۱۰، ۲۰ و ۴۰ درصد): تغییر خروجی (واحد مدل) وقتی ورودی‌هایی که هر روش مهم‌تر می‌داند حذف می‌شوند. بزرگ‌تر بهتر است.}
\label{tab:interp-del}
\begin{tabular}{cccccc}
\toprule
خروجی & حذف تصادفی & \lr{rollout} & \lr{attention flow} & \lr{gradient-weighted rollout} & \lr{gradient × input} \\
\midrule
\lr{E} & \lr{۰٫۰۲۱} & \lr{۰٫۰۹۸} & \lr{۰٫۰۹۳} & \lr{۰٫۲۰۵} & \lr{۰٫۰۹۲} \\
\lr{K} & \lr{۰٫۰۳۷} & \lr{۰٫۲۱۲} & \lr{۰٫۱۱۸} & \lr{۰٫۲۱۳} & \lr{۰٫۲۱۱} \\
\bottomrule
\end{tabular}
\end{table}


\paragraph{کدام نقشه‌ها واقعاً از یادگیری می‌آیند؟} جدول~\ref{tab:interp-sanity} یک هشدار دارد. \lr{plain rollout} \lr{K} برای مدل به‌هم‌ریخته هم شکل مشابهی دارد (همبستگی اسپیرمن \lr{۰٫۴۷}). دلیلش ساختار مدل است: \lr{residual connection} و ردپای \lr{stem} همیشه بیشترین سهم را به \lr{token} خودِ رخداد و همسایه‌هایش می‌دهند، چه مدل آموزش دیده باشد چه نه. برای \lr{E}، \lr{gradient × input} کمترین وابستگی به یادگیری را دارد (\lr{۰٫۲۷}). \lr{gradient-weighted rollout} در هر دو خروجی نزدیک صفر است.

نتیجه‌ی این دو آزمون با هم: \lr{gradient-weighted rollout} هم ورودی‌های مهم را بهتر پیدا می‌کند و هم کمتر از ساختار مدل اثر می‌گیرد. برای همین روش اصلی ادعاهای این گزارش است.

\begin{table}[htbp]
\centering
\small
\caption{\lr{weight-shuffle check}: \lr{Spearman correlation} بین نقشه‌ی مدل آموزش‌دیده و نقشه‌ی مدلی با ماتریس‌های وزنِ به‌هم‌ریخته (میانگین روی ۱۲ شرکت‌کننده). نزدیک صفر بهتر است: یعنی نقشه به آموخته‌های مدل وابسته است.}
\label{tab:interp-sanity}
\begin{tabular}{ccccc}
\toprule
خروجی & \lr{rollout} & \lr{attention flow} & \lr{gradient-weighted rollout} & \lr{gradient × input} \\
\midrule
\lr{E} & \lr{۰٫۰۳۶} & \lr{۰٫۰۷۵} & \lr{۰٫۰۱۲} & \lr{۰٫۲۷۱} \\
\lr{K} & \lr{۰٫۴۷۳} & \lr{−۰٫۰۰۸} & \lr{۰٫۰۷۰} & \lr{۰٫۰۷۶} \\
\bottomrule
\end{tabular}
\end{table}

\begin{table}[htbp]
\centering
\small
\caption{اجزای دقیق خروجی در پنجره‌های رخداد و بهترین روش تفسیر برای \lr{E} در هر شرکت‌کننده.}
\label{tab:interp-parts}
\begin{tabular}{ccccc}
\toprule
شرکت‌کننده & بازه‌ی \lr{history gain} & $|F_{eeg}|/K$ (\%) & بیشینه‌ی $|g|$ & بهترین روش در \lr{deletion check} (\lr{E}) \\
\midrule
\lr{P1} & \lr{۱٫۰۹}--\lr{۱٫۲۰} & \lr{۰٫۶۹} & \lr{۰٫۰۱۵} & \lr{gradient-weighted rollout} \\
\lr{P2} & \lr{۱٫۰۶}--\lr{۱٫۱۳} & \lr{۰٫۱۵} & \lr{۰٫۰۱۴} & \lr{gradient-weighted rollout} \\
\lr{P3} & \lr{۱٫۱۰}--\lr{۱٫۱۸} & \lr{۱٫۰۱} & \lr{۰٫۰۱۷} & \lr{gradient-weighted rollout} \\
\lr{P4} & \lr{۱٫۱۰}--\lr{۱٫۱۹} & \lr{۰٫۶۰} & \lr{۰٫۰۲۶} & \lr{gradient-weighted rollout} \\
\lr{P5} & \lr{۱٫۰۴}--\lr{۱٫۱۱} & \lr{۰٫۵۷} & \lr{۰٫۰۲۶} & \lr{gradient-weighted rollout} \\
\lr{P6} & \lr{۱٫۰۸}--\lr{۱٫۱۵} & \lr{۱٫۱۴} & \lr{۰٫۰۲۴} & \lr{gradient-weighted rollout} \\
\lr{P7} & \lr{۱٫۱۱}--\lr{۱٫۱۸} & \lr{۰٫۹۱} & \lr{۰٫۰۱۸} & \lr{gradient-weighted rollout} \\
\lr{P8} & \lr{۱٫۰۸}--\lr{۱٫۱۸} & \lr{۲٫۲۲} & \lr{۰٫۰۴۲} & \lr{gradient-weighted rollout} \\
\lr{P9} & \lr{۱٫۰۸}--\lr{۱٫۱۸} & \lr{۰٫۶۵} & \lr{۰٫۰۱۲} & \lr{gradient-weighted rollout} \\
\lr{P10} & \lr{۱٫۰۹}--\lr{۱٫۲۶} & \lr{۰٫۸۳} & \lr{۰٫۰۱۰} & \lr{gradient-weighted rollout} \\
\lr{P11} & \lr{۱٫۰۸}--\lr{۱٫۱۴} & \lr{۱٫۰۹} & \lr{۰٫۰۱۵} & \lr{gradient-weighted rollout} \\
\lr{P12} & \lr{۱٫۰۹}--\lr{۱٫۱۷} & \lr{۱٫۰۶} & \lr{۰٫۰۲۱} & \lr{gradient-weighted rollout} \\
\bottomrule
\end{tabular}
\end{table}


\subsection{نتایج \lr{P1} و \lr{P2}: هر شکل چه می‌گوید}
\label{sec:interp-figs}

همه‌ی شکل‌های این بخش برای رخداد «جدا شدن» (\lr{lift\_off}) و میانگین روی همه‌ی بلندکردن‌های سری‌های آزمون هستند. شکل‌ها بر اساس نوعشان دسته‌بندی شده‌اند و از هر نوع، شکل \lr{P1} و \lr{P2} کنار هم آمده تا بتوان دو نفر را مقایسه کرد. برای هر نوع شکل سه چیز توضیح داده می‌شود: چه چیزی رسم شده، چطور خوانده می‌شود و چه نتیجه‌ای از آن می‌گیریم.

\subsubsection*{۱. شکل زمانی \lr{K}: مدل سینماتیک به کدام لحظه‌ها نگاه می‌کند؟}

\textbf{چه رسم شده.} پنج پنل بالا یکی برای هر عضله است. محور افقی زمان نسبت به لحظه‌ی جدا شدن است، از ۴ ثانیه قبل تا ۴ ثانیه بعد (کل پنجره‌ی ۸ ثانیه‌ای مدل). محور عمودی سهم هر \lr{token} ورودی در پیش‌بینی \lr{K} در لحظه‌ی صفر است. آبی \lr{rollout} و سبز \lr{gradient-weighted rollout} است و نوار کم‌رنگ دور هر منحنی \lr{SEM} روی بلندکردن‌هاست. سهم خودِ \lr{token} رخداد و همسایه‌هایش بالای هر پنل به‌صورت عدد نوشته شده (\mbox{«\lr{own token: 14\%}»}) و منحنی سهم بقیه‌ی \lr{tokens} است. پنل پایین \lr{attention flow} منهای جریان توجه یکنواخت است؛ صفر یعنی «مثل مدلی که به همه‌جا یکسان نگاه می‌کند» و منفی یعنی مدل راه آن لحظه را تنگ کرده است.

\textbf{چطور بخوانیم.} جایی که منحنی بالاتر است، مدل برای پیش‌بینی فعالیت عضله در لحظه‌ی جدا شدن بیشتر از اطلاعات آن لحظه استفاده کرده است. اگر منحنی صاف بود، یعنی همه‌ی لحظه‌ها به یک اندازه مهم بودند و مدل جای خاصی را ترجیح نمی‌داد.

\textbf{چه می‌فهمیم.} در هر دو نفر الگو یکی است و سه بخش دارد:
\begin{itemize}
\item \textbf{پیش از رخداد (۴ ثانیه تا صفر):} سهم بالاست. در \lr{P2} از حدود ۲ ثانیه پیش از جدا شدن بالا می‌رود و درست پیش از صفر به اوج می‌رسد. این همان بازه‌ای است که دست به جسم می‌رسد، آن را می‌گیرد و نیروی بار را بالا می‌برد. پس \lr{K} برای پیش‌بینی فعالیت عضله در لحظه‌ی جدا شدن، به حرکت و نیروی پیش از آن نگاه می‌کند.
\item \textbf{۰ تا حدود \lr{۲٫۵} ثانیه بعد:} سهم کمترین است و جریان هم بیشترین افت را دارد (در \lr{P2} تا \lr{−۰٫۴۵}). این بازه آغاز نگه‌داشتن است، اما آن لحظه‌ها هنوز بعد از رخدادند و برای توضیح فعالیت در لحظه‌ی جدا شدن اطلاعات تازه‌ای ندارند.
\item \textbf{حدود ۳ تا ۴ ثانیه بعد:} سهم دوباره بالا می‌رود. بخشی از این افزایش اثر لبه‌ی پنجره است که در همه‌ی افراد دیده می‌شود (بخش~\ref{sec:consistency}) و نباید آن را به رویداد حرکتی نسبت داد.
\end{itemize}
منحنی آبی در پنج پنل تقریباً یکی است، چون \lr{rollout} زمانی به عضله بستگی ندارد. منحنی سبز برای هر عضله قله‌های خودش را دارد؛ مثلاً در \lr{P1} \lr{ED} قله‌ی بلندی حدود \lr{۰٫۸} ثانیه پیش از رخداد دارد. یعنی گرادیان نشان می‌دهد هر عضله به لحظه‌ی کمی متفاوتی از حرکت حساس است.

\fig{P1_time_K_lift_off.png}{1.0}{\lr{P1}، خروجی \lr{K}: سهم هر \lr{token} زمانی در اطراف جدا شدن جسم (بالا) و \lr{attention flow} منهای \lr{uniform-attention flow} (پایین).}
\fig{P2_time_K_lift_off.png}{1.0}{\lr{P2}، خروجی \lr{K}: سهم زمانی و \lr{attention flow} منهای خط پایه. همان الگوی \lr{P1}: سهم بالا پیش از رخداد، کمترین سهم در ۰ تا \lr{۲٫۵} ثانیه بعد.}

\subsubsection*{۲. شکل زمانی \lr{E}: مدل \lr{EEG} به کدام لحظه‌ها نگاه می‌کند؟}

\textbf{چه رسم شده.} مثل شکل قبل، اما برای خروجی \lr{E} که فقط \lr{EEG} می‌بیند.

\textbf{چه می‌فهمیم.} سه تفاوت با \lr{K} دارد:
\begin{itemize}
\item \textbf{سهم خودِ لحظه بیشتر است:} در \lr{P1} حدود ۲۶ درصد در برابر ۱۴ درصد برای \lr{K}. یعنی \lr{E} بیشتر به توان باندهای \lr{EEG} در همان لحظه تکیه می‌کند و کمتر به گذشته.
\item \textbf{منحنی صاف‌تر است:} سهم لحظه‌های دیگر تقریباً یکسان است و ساختار زمانی روشنی مثل \lr{K} ندارد. در \lr{P1} سبز کمی پیش از صفر بالا می‌رود که با زمان‌بندی \lr{ERD} پیش از حرکت سازگار است. در \lr{P2} حتی این هم دیده نمی‌شود.
\item \textbf{جهش‌های تیز در اواخر پنجره:} در \lr{P1} یک قله‌ی تیز حدود \lr{۳٫۲} ثانیه بعد از رخداد هست و در \lr{P2} جهش در لبه‌ی پنجره. چون این قله‌ها در همه‌ی افراد در یک زمان نیستند (بخش~\ref{sec:consistency})، نمی‌توان آن‌ها را به یک رویداد حرکتی مشخص نسبت داد.
\end{itemize}
پنل جریان برای \lr{E} تقریباً صاف و همه‌جا حدود \lr{−۰٫۲} است. یعنی جریان در مسیر \lr{EEG} درباره‌ی زمان تقریباً چیزی نمی‌گوید.

\fig{P1_time_E_lift_off.png}{1.0}{\lr{P1}، خروجی \lr{E}: سهم زمانی و \lr{attention flow} منهای خط پایه. سهم خودِ لحظه ۲۶ درصد است، بیشتر از \lr{K}.}
\fig{P2_time_E_lift_off.png}{1.0}{\lr{P2}، خروجی \lr{E}: سهم زمانی و \lr{attention flow} منهای خط پایه. منحنی تقریباً صاف است و فقط در لبه‌ی پنجره جهش دارد.}

\subsubsection*{۳. نقشه‌ی سر: \lr{E} از کدام الکترودها اطلاعات می‌گیرد؟}

\textbf{چه رسم شده.} هر دایره سر از بالاست (بینی بالا، گوش‌ها در دو طرف) و نقطه‌ها جای ۳۲ الکترود. رنگ سهم هر الکترود است: هر چه تیره‌تر، سهم بیشتر. سهم یکنواخت برای هر الکترود $1/32\approx0.031$ است؛ پس رنگ‌های تیره‌تر از وسط نوار رنگ یعنی سهمی بیشتر از حالت یکنواخت. اولین دایره \lr{rollout} است (برای همه‌ی عضله‌ها یکی) و پنج دایره‌ی بعد \lr{gradient-weighted rollout} برای هر عضله. بین الکترودها درون‌یابی شده تا نقشه پیوسته دیده شود.

\textbf{چطور بخوانیم.} کانون‌های تیره جاهایی‌اند که \lr{E} بیشترین اطلاعات را از آن‌ها گرفته است. ناحیه‌ی مرکزی چپ (\lr{C3}، \lr{FC5}، \lr{CP5}) بالای قشر حرکتی دست راست است و انتظار داریم آنجا مهم باشد. پیشانی (\lr{Fp1}، \lr{Fp2}، \lr{F7}، \lr{F8}) به حرکت چشم حساس است، و پس‌سری و گیجگاهی (\lr{O1}، \lr{O2}، \lr{PO9}، \lr{PO10}، \lr{TP9}، \lr{TP10}) به پردازش بینایی و فعالیت عضلات گردن.

\textbf{چه می‌فهمیم.}
\begin{itemize}
\item در \lr{P1} قوی‌ترین کانون روی مرکزی چپ است؛ یعنی همان جایی که از نظر فیزیولوژی انتظار داریم. کانون کوچک‌تری روی مرکزی راست (\lr{C4}) و کانون قوی دیگری روی پس‌سری هم هست.
\item در \lr{P2} کانون مرکزی چپ ضعیف‌تر است و بیشترین سهم روی لبه‌های سر، به‌خصوص پیشانی راست و پس‌سری، است. این نشانه‌ی احتمال \lr{artifact} است.
\item پنج نقشه‌ی عضله‌ها در هر نفر تقریباً یکسان‌اند. پس \lr{E} برای همه‌ی عضله‌ها از یک جای سر می‌خواند و نقشه‌ی خاص هر عضله نمی‌سازد؛ این با خواندن «مرحله‌ی کلی حرکت» سازگار است، نه فرمان جداگانه برای هر عضله.
\item در \lr{P1} نقشه‌ی \lr{attention flow} رسم نشده، چون سهم همه‌ی الکترودها بین \lr{۳٫۰۹} و \lr{۳٫۱۵} درصد است و نقشه چیزی جز رنگ یکسان نشان نمی‌داد. در \lr{P2} رسم شده، اما نوار رنگش فقط از \lr{۰٫۰۲۹} تا \lr{۰٫۰۳۲} است؛ تفاوت‌ها بسیار کوچک‌اند.
\end{itemize}

\fig{P1_scalp_E_lift_off.png}{1.0}{\lr{P1}، خروجی \lr{E}: سهم هر الکترود در \lr{rollout} و در \lr{gradient-weighted rollout} برای هر عضله. کانون اصلی روی مرکزی چپ (\lr{C3}). نقشه‌ی \lr{attention flow} رسم نشده، چون برای همه‌ی الکترودها تقریباً برابر است.}
\fig{P2_scalp_E_lift_off.png}{1.0}{\lr{P2}، خروجی \lr{E}: سهم هر الکترود. کانون مرکزی چپ ضعیف‌تر و سهم لبه‌های سر بیشتر است. جریان الکترودها کمی پراکندگی دارد (\lr{۱۱٫۵} درصد)، پس نقشه‌اش با رنگ از کمینه تا بیشینه رسم شده است.}

\subsubsection*{۴. توجه هر لایه‌ی \lr{electrode graph}: الکترودها از هم چه می‌خوانند؟}

\textbf{چه رسم شده.} دو ماتریس $32\times32$، یکی برای هر لایه‌ی \lr{electrode graph}، در لحظه‌ی جدا شدن و میانگین روی سرها. هر سطر الکترودی است که اطلاعات می‌گیرد و هر ستون الکترودی که اطلاعات می‌دهد. جمع هر سطر ۱ است.

\textbf{چطور بخوانیم.} یک ستون تیره که از بالا تا پایین ادامه دارد یعنی همه‌ی الکترودها از همان یک الکترود می‌خوانند؛ آن الکترود یک «مرکز پخش» است. قطر تیره یعنی هر الکترود بیشتر اطلاعات خودش را نگه می‌دارد.

\textbf{چه می‌فهمیم.} لایه‌ی اول در هر دو نفر چند مرکز پخش دارد و لایه‌ی دوم بیشتر قطری است. یعنی لایه‌ی اول اطلاعات چند الکترود خاص را به همه می‌رساند و لایه‌ی دوم آن را با اطلاعات خود هر الکترود ترکیب می‌کند. نکته‌ی مهم این است که مرکزها کدام‌اند:
\begin{itemize}
\item در \lr{P1}: \lr{C3} (قشر حرکتی چپ) و \lr{O1}/\lr{PO9} (پس‌سری). \lr{C3} در لایه‌ی دوم هم یک ستون تیره دارد.
\item در \lr{P2}: \lr{F8} و \lr{F7} (پیشانی، نزدیک چشم‌ها) و \lr{PO10}/\lr{TP10} (پشت گوش و گردن).
\end{itemize}
این شکل توضیح می‌دهد نقشه‌ی سر از کجا آمده است: نقشه‌ی \lr{P1} بیشتر از قشر حرکتی می‌آید و نقشه‌ی \lr{P2} بیشتر از الکترودهایی که به حرکت چشم و عضلات گردن حساس‌اند. پس اطلاعات \lr{E} در \lr{P2} احتمالاً بیشتر از \lr{artifact} است تا از قشر حرکتی.

\fig{P1_electrode_layers_lift_off.png}{0.8}{\lr{P1}: توجه دو لایه‌ی \lr{electrode graph} در لحظه‌ی جدا شدن (سطر = الکترود گیرنده). مرکزهای پخش در لایه‌ی اول: \lr{C3} و \lr{O1}.}
\fig{P2_electrode_layers_lift_off.png}{0.8}{\lr{P2}: توجه دو لایه‌ی \lr{electrode graph} در لحظه‌ی جدا شدن. مرکزهای پخش در لایه‌ی اول: \lr{F8}، \lr{F7} و \lr{PO10}.}

\subsubsection*{۵. \lr{muscle graph}: عضله‌ها از هم اطلاعات می‌گیرند؟}

\textbf{چه رسم شده.} چهار ماتریس $5\times5$: \lr{rollout} و \lr{gradient-weighted rollout} (\lr{relevance}) برای \lr{K} و \lr{E}. سطر عضله‌ی گیرنده و ستون عضله‌ی فرستنده است.

\textbf{چطور بخوانیم.} قطر تیره یعنی هر عضله پیش‌بینی‌اش را بیشتر از گره خودش می‌سازد. خانه‌ی تیره بیرون قطر، مثلاً سطر \lr{AD} و ستون \lr{ED}، یعنی پیش‌بینی \lr{AD} از اطلاعات گره \lr{ED} هم استفاده کرده است.

\textbf{چه می‌فهمیم.} در هر دو نفر و هر دو خروجی قطر غالب است؛ عضله‌ها بیشتر مستقل پیش‌بینی می‌شوند. در \lr{rollout} خانه‌های بیرون قطر تقریباً یکسان‌اند و الگوی خاصی ندارند. در \lr{relevance} چند خانه‌ی بیرون قطر برجسته می‌شوند؛ مثلاً در \lr{P1} \lr{AD} از \lr{ED} برای \lr{K}، و \lr{BR} از \lr{FD} برای \lr{E}. اما این رابطه‌ها از یک نفر به نفر دیگر عوض می‌شوند و بخش~\ref{sec:consistency} نشان می‌دهد در \lr{E} یکسانی‌شان بین افراد کم است. پس نباید آن‌ها را هم‌کاری ثابت عضله‌ها تعبیر کرد.

\fig{P1_muscle_graph_lift_off.png}{0.95}{\lr{P1}: \lr{muscle graph} در لحظه‌ی جدا شدن برای \lr{K} و \lr{E}، با \lr{rollout} و \lr{gradient-weighted rollout}.}
\fig{P2_muscle_graph_lift_off.png}{0.95}{\lr{P2}: \lr{muscle graph} در لحظه‌ی جدا شدن.}

\subsubsection*{۶. مقایسه با استراحت: کدام الکترود فقط هنگام حرکت مهم می‌شود؟}

\textbf{چه رسم شده.} برای هر عضله یک نقشه‌ی سر. رنگ آماره‌ی $t$ جفت‌شده است: سهم الکترود در لحظه‌ی جدا شدن منهای سهمش در استراحت بعد از همان بلندکردن. قرمز یعنی در رخداد سهم بیشتری دارد، آبی یعنی در استراحت. دایره‌های توخالی الکترودهایی‌اند که اختلافشان پس از تصحیح برای ۳۲ الکترود معنادار است.

\textbf{چطور بخوانیم.} نقشه‌ی سر (شکل نوع ۳) فقط می‌گوید \lr{E} از کجا می‌خواند؛ ممکن است همان الکترود در استراحت هم همان‌قدر مهم باشد. این شکل جواب می‌دهد کدام الکترودها \emph{مخصوص لحظه‌ی حرکت} مهم می‌شوند.

\textbf{چه می‌فهمیم.} نتیجه بین دو نفر و حتی بین عضله‌های یک نفر بسیار متفاوت است. در \lr{P1} بیشترین الکترودهای معنادار برای \lr{FD} است (پیشانی قرمز، پس‌سری آبی) و برای بقیه‌ی عضله‌ها چند الکترود پراکنده. در \lr{P2} برای همه‌ی عضله‌ها پس‌سری (\lr{O1}، \lr{O2}، \lr{P7}، \lr{P8}) قرمز و معنادار است و پیشانی راست آبی. افزایش سهم پس‌سری درست هنگام جدا شدن جسم با نگاه کردن به جسم و پردازش بینایی سازگار است، نه با فرمان حرکتی. بخش~\ref{sec:consistency} این را تأیید می‌کند: این نقشه‌ها بین افراد تقریباً شبیه نیستند (\lr{pairwise $\rho$} حدود \lr{۰٫۰۶})، پس برای نتیجه‌گیری گروهی مناسب نیستند.

\fig{P1_rest_contrast_lift_off.png}{0.95}{\lr{P1}: سهم الکترودها در \lr{E} هنگام جدا شدن منهای استراحت بعد از آن؛ آماره‌ی $t$ جفت‌شده. دایره‌ها الکترودهایی‌اند که $p$ تصحیح‌شده‌شان کمتر از \lr{۰٫۰۵} است.}
\fig{P2_rest_contrast_lift_off.png}{0.95}{\lr{P2}: مقایسه‌ی رخداد با استراحت برای \lr{E}. پس‌سری در همه‌ی عضله‌ها هنگام جدا شدن سهم بیشتری دارد.}

\subsubsection*{۷. سهم ورودی‌ها: کدام ویژگی‌ها و کدام باندها؟}

\textbf{چه رسم شده.} سه نمودار میله‌ای از \lr{gradient × input}، میانگین روی پنج عضله. چپ: سهم هر یک از ۳۲ ویژگی \lr{K} (۲۴ ویژگی سینماتیک و ۸ ویژگی تاریخچه، با نام‌های کد؛ پیشوند \lr{v\_} یعنی سرعت). وسط: سهم چهار باند \lr{EEG} در \lr{E}. راست: سهم همان چهار باند در اصلاح \lr{EEG} درون \lr{F}. سهم‌های هر نمودار جمع ۱ دارند.

\textbf{چطور بخوانیم.} این تنها شکلی است که درباره‌ی \emph{نوع} ورودی می‌گوید، نه جای آن در زمان یا روی سر. میله‌ی بلندتر یعنی تغییر کوچک آن ویژگی خروجی را بیشتر تغییر می‌دهد.

\textbf{چه می‌فهمیم.}
\begin{itemize}
\item \textbf{\lr{K} از فیزیک بلند کردن استفاده می‌کند.} در \lr{P1} بیشترین سهم با ارتفاع شست و انگشت اشاره (\lr{p\_thumb\_z}، \lr{p\_index\_z}) و سرعت نیروی بار و گرفتن (\lr{v\_F\_L\_gated}، \lr{v\_F\_G\_gated}) است. در \lr{P2} خود دو نیرو (\lr{F\_G\_gated}، \lr{F\_L\_gated}) و سرعت باز و بسته شدن دست (\lr{v\_d\_grip}) جلوترند. لحظه‌ی جدا شدن دقیقاً جایی است که نیروی بار به وزن جسم می‌رسد و دست بالا می‌رود، پس این ویژگی‌ها همان چیزی‌اند که از فیزیک انتظار داریم.
\item \textbf{تاریخچه کمک می‌کند ولی نقش اصلی ندارد.} هشت ویژگی تاریخچه (از \lr{has history} تا \lr{mean grip}) هر کدام ۱ تا ۵ درصد سهم دارند. \lr{has history} از همه بیشتر است؛ یعنی مدل بلندکردن اول سری را از بقیه جدا می‌کند.
\item \textbf{\lr{E} باند بتا را ترجیح می‌دهد.} در هر دو نفر بتا (۱۳ تا ۳۰ هرتز) اول است و گامای پایین (۳۰ تا ۴۰ هرتز) دوم. بتا با پژوهش‌های قشر حرکتی می‌خواند. سهم بالای گامای پایین به \lr{artifact} عضلات سر حساس است.
\item \textbf{اصلاح \lr{EEG} در \lr{F} همان باندهای \lr{E} را می‌خواند.} نمودار راست تقریباً تکرار نمودار وسط است. پس \lr{fusion} اطلاعات دیگری از \lr{EEG} نمی‌گیرد، فقط همان را با \lr{gate} نزدیک صفر اضافه می‌کند.
\end{itemize}

\fig{P1_inputs_lift_off.png}{0.85}{\lr{P1}: سهم ویژگی‌های سینماتیک و تاریخچه در \lr{K} و سهم باندهای \lr{EEG} در \lr{E} و در اصلاح \lr{EEG} درون \lr{F} (\lr{gradient × input}).}
\fig{P2_inputs_lift_off.png}{0.85}{\lr{P2}: سهم ورودی‌ها (\lr{gradient × input}).}

\subsubsection*{۸. \lr{deletion check}: کدام روش تفسیر را باور کنیم؟}

\textbf{چه رسم شده.} محور افقی درصد ورودی‌هایی است که حذف شده‌اند (۵، ۱۰، ۲۰ و ۴۰ درصد)، به ترتیبی که هر روش مهم‌تر می‌داند. محور عمودی تغییر خروجی مدل در اثر این حذف است. هر رنگ یک روش است و خاکستری حذف تصادفی. چپ \lr{E} و راست \lr{K}.

\textbf{چطور بخوانیم.} این شکل خود روش‌ها را می‌سنجد، نه مدل را. روشی که واقعاً ورودی‌های مهم را پیدا کرده باشد، با حذف همان ۵ درصد اول باید خروجی را بسیار تغییر دهد؛ یعنی منحنی‌اش از همان ابتدا بالا باشد. منحنی نزدیک خاکستری یعنی آن روش بهتر از انتخاب تصادفی نیست.

\textbf{چه می‌فهمیم.}
\begin{itemize}
\item \textbf{برای \lr{E}:} سبز (\lr{gradient-weighted rollout}) در هر دو نفر از همان ۵ درصد اول بسیار بالاتر از بقیه است (در \lr{P1} حدود \lr{۰٫۲۸} در برابر حدود \lr{۰٫۰۸}). پس برای \lr{EEG} فقط این روش ورودی‌های مهم را دقیق رتبه‌بندی می‌کند. در \lr{P2} سبز بعد از ۲۰ درصد کمی پایین می‌آید؛ احتمالاً با حذف ورودی‌های بیشتر، اثرهای مخالف هم حذف می‌شوند و یکدیگر را خنثی می‌کنند.
\item \textbf{برای \lr{K}:} سه روش \lr{rollout}، \lr{gradient-weighted rollout} و \lr{gradient × input} تقریباً روی هم‌اند. در \lr{P1} منحنی آن‌ها از ۵ درصد صاف است؛ یعنی تقریباً همه‌ی اثر در همان ۵ درصد اول است (خودِ لحظه‌ی رخداد و همسایه‌هایش) و حذف بیشتر چیزی اضافه نمی‌کند. \lr{attention flow} در \lr{P1} از ۵ درصد بسیار پایین شروع می‌شود، چون تقریباً به همه‌ی ورودی‌ها سهم برابر می‌دهد و نمی‌داند کدام را اول حذف کند.
\item همه‌ی روش‌ها در هر دو نفر بالاتر از حذف تصادفی‌اند؛ پس هیچ‌کدام بی‌معنا نیستند، فقط دقتشان فرق دارد.
\end{itemize}

\fig{P1_deletion.png}{0.75}{\lr{P1}: منحنی‌های \lr{deletion check}. منحنی بالاتر در درصدهای کم یعنی روش دقیق‌تر.}
\fig{P2_deletion.png}{0.75}{\lr{P2}: منحنی‌های \lr{deletion check}.}
\clearpage
